% Source: https://artofproblemsolving.com/community/c3h2030971_confusing_measure_theory
% Stable identity: aops-topic-2030971-q01
Suppose each outcome $n=1,2,3,\dots$ has the same probability $p\ge0$.
\begin{itemize}
\item If $p>0$, choose an integer $N>\dfrac{1}{p}$. The outcomes $1,2,\dots,N$ are disjoint, so the probability that one of them occurs is $Np>1$, which is impossible.
\item If $p=0$, then the total probability is $0+0+0+\cdots=0$ (every partial sum is $0$), not $1$.
\end{itemize}
In either case the probabilities cannot add up to $1$. So no probability assignment makes infinitely many outcomes $1,2,3,\dots$ equally likely. $\blacksquare$

\textbf{The error.} The argument assumes, without justification, that all the distances are equally likely. By the proof, that assumption is impossible: any genuine probability model for the distance must give different outcomes different probabilities (for example $\dfrac{1}{2^n}$ for distance $n$, which add up to $1$). Also, $\dfrac{1}{\infty}$ is not a number; the ``contradiction'' comes from the false premise, not from probability theory.
